\documentclass{ximera}

\input{../../../preamble.tex}


\definecolor{fvdnavy}{HTML}{071D43}
\definecolor{fvdcyan}{HTML}{20BFC6}
\definecolor{fvdcyanlight}{HTML}{EFFCFB}
\definecolor{fvdmagenta}{HTML}{F10D69}
\definecolor{fvdmagentalight}{HTML}{FFF1F7}
\definecolor{fvdtext}{HTML}{163044}





%=================================================
% Pedagogische omgevingen
% PDF: tcolorbox
% HTML: learning-box
%=================================================

\newenvironment{voorkennis}
{%
    \ifdefined\HCode
        \HCode{%
<section class="learning-box learning-box--prior">
<div class="learning-box__header">
<span class="learning-box__icon" aria-hidden="true"></span>
<span class="learning-box__title">Voorkennis</span>
</div>
<div class="learning-box__content">
        }%
        \begin{itemize}
    \else
       \begin{tcolorbox}[
    enhanced,
    breakable,
    colback=fvdcyanlight,
    colframe=fvdcyan,
    coltext=fvdtext,
    boxrule=0.5pt,
    leftrule=4pt,
    arc=4mm,
    outer arc=4mm,
    left=7mm,
    right=6mm,
    top=5mm,
    bottom=5mm,
    before skip=6mm,
    after skip=6mm,
    title=Voorkennis,
    coltitle=fvdcyan,
    fonttitle=\bfseries\Large,
    attach title to upper={\par\smallskip},
    boxed title style={
        empty,
        left=0mm,
        right=0mm,
        top=0mm,
        bottom=0mm
    },
    overlay={
        \draw[
            fvdcyan,
            line width=0.5pt
        ]
        ([xshift=42mm,yshift=-13mm]frame.north west)
        --
        ([xshift=-7mm,yshift=-13mm]frame.north east);
    }
]
        \begin{itemize}
    \fi
}
{%
    \end{itemize}
    \ifdefined\HCode
        \HCode{%
</div>
</section>
        }%
    \else
        \end{tcolorbox}
    \fi
}


\newenvironment{leerdoelen}
{%
    \ifdefined\HCode
        \HCode{%
<section class="learning-box learning-box--goals">
<div class="learning-box__header">
<span class="learning-box__icon" aria-hidden="true"></span>
<span class="learning-box__title">Leerdoelen</span>
</div>
<div class="learning-box__content">
        }%
        \begin{itemize}
    \else
\begin{tcolorbox}[
    enhanced,
    breakable,
    colback=fvdmagentalight,
    colframe=fvdmagenta,
    coltext=fvdtext,
    boxrule=0.5pt,
    leftrule=4pt,
    arc=4mm,
    outer arc=4mm,
    left=7mm,
    right=6mm,
    top=5mm,
    bottom=5mm,
    before skip=6mm,
    after skip=6mm,
    title=Leerdoelen,
    coltitle=fvdmagenta,
    fonttitle=\bfseries\Large,
    attach title to upper={\par\smallskip},
    boxed title style={
        empty,
        left=0mm,
        right=0mm,
        top=0mm,
        bottom=0mm
    },
    overlay={
        \draw[
            fvdmagenta,
            line width=0.5pt
        ]
        ([xshift=42mm,yshift=-13mm]frame.north west)
        --
        ([xshift=-7mm,yshift=-13mm]frame.north east);
    }
]
        \begin{itemize}
    \fi
}
{%
    \end{itemize}
    \ifdefined\HCode
        \HCode{%
</div>
</section>
        }%
    \else
        \end{tcolorbox}
    \fi
}


\addPrintStyle{../../..}

\title{Oefeningen --- Van data naar puntenwolk}
\author{Ingmar Herreman}



\begin{document}

% FVD_CHAPTER_DATA_START
% Automatisch gegenereerd uit index.tex.
% Niet handmatig aanpassen.
\fvdchapterdata{3}{Verbanden onderzoeken met data}{5}
% FVD_CHAPTER_DATA_END




\maketitle

% ============================================================
\FVDsection{Basis --- gegevens en spreidingsdiagrammen}
% ============================================================

\begin{oefening}[Gekoppelde gegevens]{\oefB \ESS}
Een bioloog meet bij zes planten de hoogte en de massa van de bovengrondse
delen.

\[
\begin{array}{c|rrrrrr}
\text{plant} &1&2&3&4&5&6\\ \hline
\text{hoogte (cm)}&12&18&23&29&34&41\\
\text{massa (g)}&4{,}1&6{,}0&7{,}4&9{,}8&10{,}7&13{,}2
\end{array}
\]

\begin{question}
Schrijf de zes gekoppelde gegevensparen op.

\begin{oplossing}
\[
(12,4{,}1),\ (18,6{,}0),\ (23,7{,}4),\ (29,9{,}8),\
(34,10{,}7),\ (41,13{,}2).
\]
\end{oplossing}
\end{question}

\begin{question}
Wat betekent het punt $(29,9{,}8)$ in deze context?

\begin{oplossing}
Plant 4 heeft een hoogte van $29\,\mathrm{cm}$ en een bovengrondse
massa van $9{,}8\,\mathrm{g}$.
\end{oplossing}
\end{question}

\begin{question}
Welke grootheid zou je hier het meest natuurlijk op de horizontale as
plaatsen?

\begin{oplossing}
De hoogte. In deze analyse onderzoeken we hoe de massa samenhangt met
de hoogte. Een andere keuze is wiskundig mogelijk, maar moet duidelijk
worden gemotiveerd.
\end{oplossing}
\end{question}
\end{oefening}


\begin{oefening}[Richting herkennen]{\oefB \ESS}
Duid voor elke situatie aan welk globaal verband je verwacht.
Dit is een \emph{verwachting}; echte gegevens kunnen afwijken.

\begin{question}
Buitentemperatuur en dagelijks gasverbruik voor verwarming:

\begin{multipleChoice}
  \choice{positief}
  \choice[correct]{negatief}
  \choice{geen mogelijk verband}
\end{multipleChoice}
\end{question}

\begin{question}
Lengte en massa van gelijkvormige metalen staven uit hetzelfde materiaal:

\begin{multipleChoice}
  \choice[correct]{positief}
  \choice{negatief}
  \choice{geen mogelijk verband}
\end{multipleChoice}
\end{question}

\begin{question}
Hoogte boven zeeniveau en luchtdruk in een geschikt meetbereik:

\begin{multipleChoice}
  \choice{positief}
  \choice[correct]{negatief}
  \choice{geen mogelijk verband}
\end{multipleChoice}
\end{question}
\end{oefening}


\begin{oefening}[Tabel naar puntenwolk]{\oefB \ESS}
Bij een proef met een veer worden de volgende gegevens verzameld.

\[
\begin{array}{c|rrrrrr}
m\;(\mathrm{g})&50&100&150&200&250&300\\ \hline
u\;(\mathrm{cm})&1{,}2&2{,}5&3{,}7&5{,}0&6{,}1&7{,}5
\end{array}
\]

\begin{question}
Maak met ICT een spreidingsdiagram met de massa $m$ op de horizontale
as en de uitrekking $u$ op de verticale as.
\end{question}

\begin{question}
Beschrijf richting, vorm en sterkte van het zichtbare verband.

\begin{oplossing}
De puntenwolk vertoont een sterk positief en ongeveer lineair verband:
grotere massa's gaan in deze dataset samen met grotere uitrekkingen en
de punten liggen dicht rond een stijgende rechte band.
\end{oplossing}
\end{question}

\begin{question}
Mag je op basis van deze grafiek alleen besluiten dat de massa de
uitrekking veroorzaakt?

\begin{oplossing}
Nee. De grafiek toont samenhang. Voor een causale conclusie is ook de
experimentele opzet en fysische kennis nodig.
\end{oplossing}
\end{question}
\end{oefening}


% ============================================================
\FVDsection{Verdieping --- puntenwolken analyseren}
% ============================================================

\begin{oefening}[Temperatuur en energieverbruik]{\oefU \ESS}
Een school onderzoekt tijdens tien winterdagen de gemiddelde
buitentemperatuur en het gasverbruik.

\[
\begin{array}{c|rrrrrrrrrr}
T\;({}^{\circ}\mathrm C)&-2&0&2&4&5&7&8&10&11&13\\ \hline
G\;(\mathrm{kWh})&820&790&735&690&655&610&590&535&520&470
\end{array}
\]

\begin{question}
Maak een spreidingsdiagram.
\end{question}

\begin{question}
Analyseer de puntenwolk met de vier aandachtspunten:
richting, vorm, sterkte en bijzonderheden.

\begin{oplossing}
Er is een sterk negatief en ongeveer lineair verband. Bij hogere
buitentemperaturen ligt het gasverbruik doorgaans lager. De punten
liggen vrij dicht rond een dalende rechte band en er is in deze kleine
dataset geen uitgesproken uitschieter zichtbaar.
\end{oplossing}
\end{question}

\begin{question}
Formuleer een contextgebonden conclusie zonder meer te beweren dan de
gegevens toelaten.

\begin{oplossing}
Tijdens deze tien winterdagen gaat een hogere gemiddelde
buitentemperatuur doorgaans samen met een lager gasverbruik van de
school. De gegevens suggereren een sterk negatief, ongeveer lineair
verband.
\end{oplossing}
\end{question}
\end{oefening}


\begin{oefening}[Een verband dat niet lineair is]{\oefU \ESS}
Een sensor wordt getest. Voor verschillende waarden van $x$ wordt
de respons $y$ gemeten.

\[
\begin{array}{c|rrrrrrrrr}
x&-4&-3&-2&-1&0&1&2&3&4\\ \hline
y&16{,}2&9{,}1&4{,}2&1{,}0&0{,}2&1{,}1&3{,}9&9{,}2&15{,}8
\end{array}
\]

\begin{question}
Maak een spreidingsdiagram.
\end{question}

\begin{question}
Is er een duidelijk verband tussen $x$ en $y$?

\begin{oplossing}
Ja. De punten volgen duidelijk een U-vormig patroon.
\end{oplossing}
\end{question}

\begin{question}
Is het verband ongeveer lineair?

\begin{oplossing}
Nee. Een rechte beschrijft het globale patroon niet goed.
\end{oplossing}
\end{question}

\begin{question}
Waarom is deze oefening belangrijk voordat we correlatie bestuderen?

\begin{oplossing}
Omdat een dataset een sterk en duidelijk verband kan vertonen zonder
dat dit verband lineair is. Een maat voor lineaire samenhang moet dus
altijd samen met de puntenwolk geïnterpreteerd worden.
\end{oplossing}
\end{question}
\end{oefening}


\begin{oefening}[De verdachte meting]{\oefU \ESS}
Bij een experiment worden de volgende waarden gemeten.

\[
\begin{array}{c|rrrrrrrr}
x&1&2&3&4&5&6&7&8\\ \hline
y&2{,}1&4{,}0&6{,}2&8{,}1&10{,}0&24{,}5&14{,}2&16{,}1
\end{array}
\]

\begin{question}
Maak een spreidingsdiagram en beschrijf het globale patroon.

\begin{oplossing}
De meeste punten volgen een sterk positief en ongeveer lineair patroon.
Het punt bij $x=6$ ligt opvallend ver boven dit patroon.
\end{oplossing}
\end{question}

\begin{question}
Welke waarneming is een mogelijke uitschieter?

\begin{oplossing}
$(6,24{,}5)$.
\end{oplossing}
\end{question}

\begin{question}
Een leerling verwijdert het punt onmiddellijk. Beoordeel die werkwijze.

\begin{oplossing}
Dat is niet verantwoord zonder onderzoek. Het kan om een invoer- of
meetfout gaan, maar ook om een correcte uitzonderlijke waarneming.
Eerst moet de oorzaak worden nagegaan.
\end{oplossing}
\end{question}
\end{oefening}


\begin{oefening}[Twee groepen in één puntenwolk]{\oefU}
Een onderzoeker meet de massa en de lengte van onderdelen die afkomstig
zijn van twee verschillende materialen. In de puntenwolk ontstaan twee
duidelijk gescheiden groepen.

\begin{question}
Hoe noemen we zulke groepen?

\begin{oplossing}
Clusters.
\end{oplossing}
\end{question}

\begin{question}
Waarom is het gevaarlijk om alleen één globale trend door alle punten
te beschrijven?

\begin{oplossing}
Omdat de twee materialen mogelijk elk een ander verband vertonen.
Een globale trend kan de groepsstructuur verbergen en daardoor een
misleidende samenvatting geven.
\end{oplossing}
\end{question}
\end{oefening}


% ============================================================
\FVDsection{Geïntegreerd analyseren}
% ============================================================

\begin{oefening}[Zonnepaneel]{\oefU \ESS}
Tijdens een practicum wordt op verschillende momenten de instraling
op een zonnepaneel en het geleverde elektrische vermogen gemeten.

\[
\begin{array}{c|rrrrrrrr}
I\;(\mathrm{W/m^2})&180&260&340&430&510&620&730&840\\ \hline
P\;(\mathrm{W})&31&43&58&70&86&101&119&136
\end{array}
\]

Voer een volledige eerste analyse uit.

\begin{question}
Benoem de twee numerieke grootheden en hun eenheden.

\begin{oplossing}
Instraling $I$ in $\mathrm{W/m^2}$ en elektrisch vermogen $P$ in
$\mathrm W$.
\end{oplossing}
\end{question}

\begin{question}
Maak met ICT een spreidingsdiagram.
\end{question}

\begin{question}
Beschrijf richting, vorm, sterkte en bijzonderheden.

\begin{oplossing}
De puntenwolk vertoont een sterk positief en ongeveer lineair verband.
Grotere instraling gaat in deze dataset samen met groter geleverd
vermogen. De punten liggen dicht rond een stijgende rechte band en er
is geen duidelijke uitschieter.
\end{oplossing}
\end{question}

\begin{question}
Formuleer een voorzichtige conclusie.

\begin{oplossing}
Binnen het onderzochte meetbereik gaat een grotere instraling doorgaans
samen met een groter elektrisch vermogen van het zonnepaneel. De
gegevens suggereren een sterk positief en ongeveer lineair verband.
\end{oplossing}
\end{question}

\begin{question}
Welke vervolgvraag ligt voor de hand als we het verband nauwkeuriger
willen beschrijven?

\begin{oplossing}
We kunnen onderzoeken hoe sterk het lineaire verband is en vervolgens
een lineair model bepalen. Dat leidt naar correlatie en lineaire
regressie.
\end{oplossing}
\end{question}
\end{oefening}


% ============================================================
\FVDsection{Gevorderd --- kritisch onderzoeken}
% ============================================================

\begin{oefening}[Kan één punt een verhaal veranderen?]{\oefG}
Een dataset bevat aanvankelijk de punten

\[
(2,8),\ (3,7),\ (4,6),\ (5,5),\ (6,4),\ (7,3).
\]

Daar wordt één extra punt $(20,20)$ aan toegevoegd.

\begin{question}
Maak een spreidingsdiagram van de eerste zes punten en beschrijf het verband.
\end{question}

\begin{question}
Voeg daarna $(20,20)$ toe. Hoe verandert het globale beeld?
\end{question}

\begin{question}
Leg uit waarom een analyse nooit uitsluitend op een automatisch
berekend getal mag steunen.

\begin{oplossing}
Een uitzonderlijk punt kan de numerieke samenvatting sterk beïnvloeden.
De puntenwolk toont de structuur van de gegevens en maakt zulke
bijzonderheden zichtbaar. Grafische en numerieke informatie moeten
daarom samen geïnterpreteerd worden.
\end{oplossing}
\end{question}
\end{oefening}


\begin{oefening}[Ontwerp zelf een dataset]{\oefG}
Maak zelf drie kleine datasets met telkens minstens acht punten:

\begin{enumerate}
  \item een sterk positief ongeveer lineair verband;
  \item een sterk negatief ongeveer lineair verband;
  \item een sterk maar duidelijk niet-lineair verband.
\end{enumerate}

Maak voor elke dataset een spreidingsdiagram en schrijf er een korte
analyse bij volgens richting, vorm, sterkte en bijzonderheden.

\begin{oplossing}
Er zijn veel correcte antwoorden mogelijk. De grafieken en de
beschrijvingen moeten onderling consistent zijn. Bij de derde dataset
moet duidelijk zijn dat ``sterk verband'' niet hetzelfde betekent als
``sterk lineair verband''.
\end{oplossing}
\end{oefening}


% ============================================================
\FVDsection{Essentiële oefeningen}
% ============================================================

Voor dit hoofdstuk behoren de volgende oefeningen tot het minimumtraject:

\begin{itemize}
  \item \textbf{Gekoppelde gegevens} --- betekenis van gegevensparen;
  \item \textbf{Richting herkennen} --- positief en negatief verband;
  \item \textbf{Tabel naar puntenwolk} --- een spreidingsdiagram maken en lezen;
  \item \textbf{Temperatuur en energieverbruik} --- een nieuwe puntenwolk analyseren;
  \item \textbf{Een verband dat niet lineair is} --- lineair en niet-lineair onderscheiden;
  \item \textbf{De verdachte meting} --- uitschieters herkennen en kritisch behandelen;
  \item \textbf{Zonnepaneel} --- geïntegreerde analyse in een STEM-context.
\end{itemize}

De eerste oefeningen trainen noodzakelijke basisvaardigheden.
De \oefU-oefeningen met \ESS zijn eveneens essentieel: zonder het
zelfstandig analyseren van nieuwe gegevens wordt het vereiste
beheersingsniveau niet bereikt.

\end{document}
